\medskip

Dans la figure ci-dessous, ABCDEFGH est un parallélépipède rectangle tel que 

AB $= 5$, AD $= 3$ et AE $= 2$.

L'espace est muni d'un repère orthonormé d'origine A dans lequel les points B, D et E ont respectivement pour coordonnées (5~;~0~;~0), (0~;~3~;~0) et (0~;~0~;~2).

Le repère est donc:
$\left (\text{A}\;; \dfrac{1}{5}\vectt{AB},\; \dfrac{1}{3}\vectt{AD},\; \dfrac{1}{2}\vectt{AE} \right )$.

\begin{center}
\psset{unit=1cm}
\begin{pspicture}(9,5)

\psframe(0.2,0.2)(7.2,3)
\psline(7.2,0.2)(8.4,1.6)(8.4,4.4)(7.2,3)
\psline(8.4,4.4)(1.6,4.4)(0.2,3)
\psline[linestyle=dashed](0.2,0.2)(1.6,1.6)(8.4,1.6)
\psline[linestyle=dashed](1.6,1.6)(1.6,4.4)
\pspolygon[linestyle=dotted,linewidth=1.5pt](0.2,0.2)(8.4,1.6)(3,4.4)
\uput[dl](0.2,0.2){A} \uput[dr](7,0.2){B} \uput[r](8.4,1.6){C}
\uput[ur](1.6,1.6){D} \uput[l](0.2,3){E} \uput[ul](7.2,3){F}
\uput[ur](8.4,4.4){G} \uput[ul](1.6,4.4){H} \uput[u](3,4.4){M}
\end{pspicture}
\end{center}

\begin{enumerate}
\item 
	\begin{enumerate}
		\item 
$\vectt{AH}= \vectt{AD} + \vectt{AE}$ donc le point H a pour coordonnées $(0\;; 3\;; 2)$.
		
$\vectt{AG}=  \vectt{AB} + \vectt{AD} + \vectt{AE}$ donc le point G a pour coordonnées $(5\;; 3\;; 2)$.		

		\item
La droite (GH) a pour vecteur directeur $\vectt{HG}$ de coordonnées $(5\;; 0\;; 0)$.

De plus, elle passe par le point G donc elle a pour représentation paramétrique:

$\left \{
\begin{array}{l !{=} l}
x & 5 +5t\\
y & 3\\
z & 2
\end{array}		
\right . \text{ avec } t\in\R$.		
	\end{enumerate}	
	
\item Soit M un point du segment [GH] tel que $\vect{\text{HM}} =k\vect{\text{HG}}$ avec $k$ un nombre réel de l'intervalle [0~;~1].
	\begin{enumerate}
		\item 
$\vectt{HM} = k \vectt{HG}
\iff
\left \{
\begin{array}{l !{=} l}
x_{\text M} -0& 5k\\
y_{\text M} -3& 0\\
z_{\text M} -2& 0
\end{array}		
\right .		
\iff
\left \{
\begin{array}{l !{=} l}
x_{\text M} & 5k\\
y_{\text M} & 3\\
z_{\text M} & 2
\end{array}		
\right .	$
		
Donc  les coordonnées de M sont $(5k~;~3~;~2)$.	
		
		\item 
Les coordonnées de $\vectt{AM}$ sont celles de M donc: $(5k~;~3~;~2)$.	
		
Les coordonnées de C sont $(5\;; 3\;; 0)$, donc celles de $\vectt{CM}$ sont $(5	k-5\;; 3-3\;; 2-0)$ soit $(5k-5\;; 0\;; 2)$.

$\vect{\text{AM}} \cdot \vect{\text{CM}}
= 5k \times (5k-5) + 3 \times 0 + 2\times 2
 = 25k^2  - 25k + 4$
	
		\item 
Le triangle AMC est rectangle en M si et seulement si $\vect{\text{AM}} \perp \vect{\text{CM}}$, c'est-à-dire $\vect{\text{AM}} \cdot \vect{\text{CM}}$ ou encore $25k^2-25k+4=0$.		
		
On résout cette équation.
$\Delta = (-25)^2-4\times 25\times 4 = 225 = 15^{2}$

L'équation admet deux solutions 
$k'=\dfrac{25+15}{2\times 25} = \dfrac{40}{50}=\dfrac{4}{5}$
et
$k''=\dfrac{25-15}{2\times 25} = \dfrac{10}{50}=\dfrac{1}{5}$

Donc pour $k=\dfrac{1}{5}$ ou $k=\dfrac{4}{5}$, le triangle AMC est rectangle en M.
		
	\end{enumerate}
\end{enumerate}	

Dans toute la suite de l'exercice, on considère que le point M a pour coordonnées (1~;~3~;~2).

On admet que le triangle AMC est rectangle en M.



\begin{enumerate}[resume]
\item On considère le point K de coordonnées (1~;~3~;~0).
	\begin{enumerate}
		\item
Le plan (ACD) a pour vecteurs directeurs $\vectt{AB}$ et $\vectt{AD}$, donc il a pour équation cartésienne $z=0$.
		
		\item 
$z_{\text K}=0$ donc le point K appartient au plan (ACD).

$\vectt{MK}$ a pour coordonnées $(0\;; 0\;; -2)$, donc \\
$\bullet~~\vectt{MK}\cdot\vectt{AB} = 0$ donc $\vectt{MK}\perp\vectt{AB}$;\\
$\bullet~~\vectt{MK}\cdot\vectt{AC} = 0$ donc $\vectt{MK}\perp\vectt{AC}$.

On en déduit que $\vectt{MK}$ est orthogonal au plan (ACD), et donc que K est le projeté orthogonal du point M sur le plan (ACD).
	
		\item Le volume du tétraèdre MACD est:
$\dfrac{1}{3}\times \text{Aire de ACD} \times \text{MK}$.

$\vectt{MK}$ a pour coordonnées $(0\;; 0\;; -2)$, donc MK $=2$.

Le triangle ACD est rectangle en D donc a pour aire 
$\dfrac{\text{AD}\times\text{DC}}{2}=\dfrac{ 3\times 5}{2}=\dfrac{15}{2}$.

Le volume du tétraèdre MACD est donc, en unités de volume:
$\dfrac{1}{3} \times \dfrac{15}{2} \times 2 = 5$.

\end{enumerate}	
	
\item On note P le projeté orthogonal du point D sur le plan (AMC).



Si on considère D comme sommet du tétraèdre MACD, la base est le triangle AMC, et la hauteur est DP.

Le triangle AMC est rectangle en M donc son aire est:
$\dfrac{\text{AM} \times \text{MC}}{2}$.

$\text{AM} = \ds\sqrt{(1-0)^2+(3-0)^2+(2-0)^2} = \ds\sqrt{14}$

$\text{CM} = \ds\sqrt{(1-5)^2+(3-3)^2+(2-0)^2}= \ds\sqrt{20}$

L'aire de AMC vaut: $\dfrac{\ds\sqrt{14}\times \sqrt{20}}{2}=\ds\sqrt{70}$.

Le tétraèdre MACD a donc pour volume: $\dfrac{1}{3}\times\ds\sqrt{70}\times\text{DP}$ soit: $\dfrac{\text{DP} \ds\sqrt{70}}{3}$.

On sait que ce volume vaut 5, donc: 
$\dfrac{\text{DP} \ds\sqrt{70}}{3} = 5$ donc:
$\text{DP}= \dfrac{15}{\ds\sqrt{70}} \approx 1,8$.
\end{enumerate}
